% Generated draft: requires mathematical and source review.
\documentclass[11pt]{article}
\usepackage[margin=25mm]{geometry}
\usepackage{fontspec}
\setmainfont{Latin Modern Roman}
\usepackage{amsmath,amssymb,mathtools,unicode-math}
\setmathfont{Latin Modern Math}
\usepackage{enumitem,needspace,xurl,hyperref}
\setlength{\parindent}{0pt}
\setlength{\parskip}{5pt}
\setlength{\emergencystretch}{3em}
\newcommand{\N}{\mathbb{N}}
\newcommand{\Z}{\mathbb{Z}}
\newcommand{\Q}{\mathbb{Q}}
\newcommand{\R}{\mathbb{R}}
\newcommand{\C}{\mathbb{C}}
\newcommand{\F}{\mathbb{F}}
\newcommand{\A}{\mathbb{A}}
\newcommand{\PP}{\mathbb{P}}
\newcommand{\E}{\mathbb{E}}
\newcommand{\eps}{\varepsilon}
\newcommand{\dd}{\,\mathrm{d}}
\begin{document}
\section*{UnsolvedMath ID: 800970}
\textbf{Source problem code:} catalogue entry
\subsection{Definitions and mathematical statement}
A closed manifold is a compact manifold without boundary. The supplied source proposes, sometimes with attribution to Berger: \emph{Let \((M,g)\) be a closed simply connected manifold such that all geodesics are closed. Then all geodesics have the same length.}

The metric signature and the convention for the length of a closed geodesic are not specified in this sentence. The following gives a precise conditional rendering under the Riemannian and least-period conventions; these conventions require human confirmation.

Assume that \(M\) is a smooth closed simply connected manifold and that \(g\) is a smooth Riemannian metric. For a unit-speed geodesic \(\gamma\colon\R\to M\), define its set of positive periods by
\[
\mathcal P(\gamma)=\{T>0:(\gamma(t+T),\dot\gamma(t+T))=(\gamma(t),\dot\gamma(t))\text{ for every }t\in\R\}.
\]
Call \(\gamma\) closed if \(\mathcal P(\gamma)\ne\varnothing\). For such a geodesic, its least positive period \(T_\gamma=\min\mathcal P(\gamma)\) is its length over one prime traversal, since \(\gamma\) has unit speed.

\textbf{Conjecture, under these conventions.} If every unit-speed geodesic on \((M,g)\) is closed, then there exists \(L>0\) such that
\[
T_\gamma=L\qquad\text{for every unit-speed geodesic }\gamma\colon\R\to M.
\]
Here the length refers to one prime traversal; repeated traversals have lengths \(kT_\gamma\) for positive integers \(k\). Constant geodesics are excluded by the unit-speed convention. The source does not explicitly state either convention.

The preceding source background separately asks for a construction or classification of all Lorentzian manifolds whose light-like geodesics are closed, citing \textnormal{[Mou-Suh2013, Suhr2013b]}. This is retained as contextual motivation, not as a Lorentzian hypothesis for the conditional conjecture above. That background supplies no dimension, compactness assumption, or precise convention for closedness. For a Lorentzian metric, a light-like geodesic has nonzero tangent satisfying \(g(\dot\gamma,\dot\gamma)=0\); its metric length does not provide the positive prime-length invariant used in the Riemannian rendering.
\subsection{Short English statement}
Under an explicitly provisional Riemannian interpretation, the conjecture says that a closed simply connected manifold on which every nonconstant geodesic is periodic has one common length for all prime geodesics. The source sometimes attributes this conjecture to Berger and labels the record open. The metric signature and prime-traversal convention need confirmation; the accompanying Lorentzian classification question is separate background.
\subsection{Sources}
\begin{itemize}
\item Local UnsolvedMath registry, permanent ID \(800970\), source problem code \texttt{catalogue entry}, title \emph{Open Problems and Questions about Geodesics --- Conjecture 51}. The supplied record lists Burns and Matveev as proposers, gives the year 2019, and labels the problem \texttt{open}; that status is supplied registry metadata, not an independently verified mathematical assessment.
\item Burns and Matveev, \emph{Open Problems and Questions about Geodesics} (2019, as identified in the supplied record), Conjecture 51: \url{https://arxiv.org/abs/1308.5417}. The supplied background records access on 2026-10-16 and extraction by \texttt{arxiv\_latex}; these are supplied provenance claims, not independently verified access or source verification.
\item The conjecture is described in the supplied statement as \emph{sometimes attributed to Berger}. No separate Berger reference is supplied.
\item The Lorentzian background cites \textnormal{[Mou-Suh2013]} and \textnormal{[Suhr2013b]}. No full bibliographic entries or links for these reference labels are supplied.
\end{itemize}
\end{document}
